\documentclass[conference]{IEEEtran}
\IEEEoverridecommandlockouts
\usepackage{cite}
\usepackage{amsmath,amssymb,amsfonts}
\usepackage{algorithmic}
\usepackage{graphicx}
\usepackage{textcomp}
\usepackage{xcolor}
\usepackage{booktabs}
\usepackage{url}
\usepackage{microtype}

\def\BibTeX{{\rm B\kern-.05em{\sc i\kern-.025em b}\kern-.08em
    T\kern-.1667em\lower.7ex\hbox{E}\kern-.125emX}}

\begin{document}

\title{AI-Driven Targeted Protein Degradation: Lead Compound Generation for Acute Myeloid Leukemia (AML) Targets via Generative PROTAC Linker Optimization}

\author{
\IEEEauthorblockN{Madhav Yalamarthi}
\IEEEauthorblockA{\textit{Department of Computer Science / Artificial Intelligence} \\
\textit{Amrita Vishwa Vidyapeetham}\\
Amritapuri, Kerala, India \\
madhav@example.com}
\and
\IEEEauthorblockN{Hareesh}
\IEEEauthorblockA{\textit{Department of Computer Science / Artificial Intelligence} \\
\textit{Amrita Vishwa Vidyapeetham}\\
Amritapuri, Kerala, India \\
hareesh@example.com}
}

\maketitle

\begin{abstract}
Targeted Protein Degradation (TPD) via Proteolysis Targeting Chimeras (PROTACs) represents a transformative therapeutic modality capable of overcoming the pharmacological limitations of classical occupancy-driven small-molecule inhibitors in Acute Myeloid Leukemia (AML). In AML, competitive inhibition of Bromodomain-containing Protein 4 (BRD4) by conventional clinical candidates (e.g., (+)-JQ1, OTX015) frequently induces compensatory homeostatic target protein accumulation and fails to abrogate non-catalytic scaffolding interactions. While heterobifunctional degraders induce catalytic, event-driven proteasomal proteolysis, rational design remains acutely constrained by empirical ``linker walking'' strategies that neglect ternary complex geometry. Here, we present an end-to-end, structure-aware computational pipeline (curate $\rightarrow$ generate $\rightarrow$ filter $\rightarrow$ model $\rightarrow$ prioritize) that couples automated ChEMBL bioactivity curation (Target ID: CHEMBL2034, sub-100~nM gate) with a Dual-Anchor Conditional Variational Autoencoder (C-VAE). The generative model conditions directly on 128-bit circular Morgan fingerprints of the warhead and E3-recruiting anchors alongside target chain length spans ($L \in [6, 22]$ heavy atoms). Over 25 epochs of training, the C-VAE achieved a 55.7\% reduction in token reconstruction loss without latent posterior collapse. In ablation benchmarking ($N = 1,000$ samples), our architecture delivered 96.8\% chemical validity and a 14.3-fold enrichment in productive lead fraction (58.7\% vs. 4.1\% for combinatorial baseline). Multiscale biophysical modeling against crystallographic references (PDB: 6BOY, 5T35) prioritized two candidate degraders: PROTAC-AML-015 ($\alpha = 8.29$, $\Delta G_{\text{ternary}} = -20.50$~kcal/mol, SAScore = 2.77) and PROTAC-AML-001 ($\alpha = 3.65$, $\Delta\text{SASA} = 542.0$~\AA$^2$, zero steric clashes). All wet-lab transitions are framed as hypotheses under a 5-phase validation cascade, providing an auditable decision framework for accelerated AML therapeutics.
\end{abstract}

\begin{IEEEkeywords}
Targeted Protein Degradation, PROTAC, Acute Myeloid Leukemia, BRD4, Cereblon, Conditional Variational Autoencoder, Ternary Complex Modeling, Thermodynamic Cooperativity.
\end{IEEEkeywords}

\section{Introduction}
Acute Myeloid Leukemia (AML) is an aggressive hematologic malignancy characterized by blocked myeloid differentiation and unrestrained self-renewal of leukemic blasts \cite{dohner2015}. Among critical epigenetic regulators, Bromodomain-containing Protein 4 (BRD4; UniProt: O60885) plays an indispensable role by binding hyperacetylated chromatin super-enhancers to sustain transcriptional elongation of major oncogenic drivers, notably \textit{MYC} and \textit{BCL2} \cite{zuber2011}.

Conventional therapeutic efforts targeting BRD4 have historically centered on competitive, occupancy-driven small-molecule bromodomain (BD1/BD2) inhibitors, including (+)-JQ1 and OTX015 (Birabresib) \cite{filippakopoulos2010}. However, clinical translation of occupancy inhibitors has been profoundly hindered by two liabilities:
\begin{enumerate}
    \item \textbf{Compensatory Target Rebound:} Pharmacological displacement of bromodomains triggers homeostatic feedback mechanisms that rapidly elevate cellular BRD4 protein abundance.
    \item \textbf{Scaffolding Persistence:} Competitive inhibitors occupy only the acetyl-lysine pocket without dismantling non-catalytic protein scaffolding networks through which BRD4 recruits P-TEFb and Mediator complexes \cite{winter2015}.
\end{enumerate}

Heterobifunctional Proteolysis Targeting Chimeras (PROTACs) offer an event-driven alternative \cite{sakamoto2001}. Structurally composed of a target-binding warhead, a chemical linker, and an E3 ubiquitin ligase recruiter (e.g., Cereblon / CRBN; UniProt: Q96SW2), a PROTAC recruits an E3 ligase into physical proximity with the protein of interest (POI). This induces formation of a transient, catalytically productive ternary complex:
\begin{equation}
\text{POI} + \text{PROTAC} + \text{E3} \rightleftharpoons [\text{Ternary Complex}] \xrightarrow{k_{\text{cat}}} \text{Ub} \rightarrow \text{Degradation}
\end{equation}

Because the degrader acts catalytically—dissociating post-polyubiquitination to initiate successive rounds of degradation—target elimination occurs at sub-stoichiometric concentrations \cite{winter2015}.

However, the primary bottleneck in PROTAC discovery is the chemical architecture of the linker. In ternary complex formation, biological degradation efficiency is driven not by individual binary affinities, but by the \textbf{thermodynamic cooperativity factor} $\alpha$:
\begin{equation}
\alpha = \frac{K_{d,\text{POI}}^{\text{binary}}}{K_{d,\text{POI}}^{\text{ternary}}} = \frac{K_{d,\text{E3}}^{\text{binary}}}{K_{d,\text{E3}}^{\text{ternary}}} = \exp\left(-\frac{\Delta\Delta G}{RT}\right)
\end{equation}

When $\alpha > 1$ (positive cooperativity), attractive inter-protein contacts stabilize the assembly ($\Delta\Delta G < 0$). When $\alpha < 1$ (negative cooperativity), steric clashes or unfavorable spatial orientations penalize assembly. Furthermore, at elevated drug concentrations, non-productive binary species outcompete ternary complex formation, causing the bell-shaped concentration-response known as the \textbf{``hook effect''} \cite{gadd2017}.

Traditional empirical optimization relies on ``linker walking''—synthesizing iterative series of alkyl or polyethylene glycol (PEG) chains. This trial-and-error process is slow, resource-intensive, and blind to protein-protein surface dynamics. To overcome this challenge, we present an end-to-end computational pipeline integrating automated ChEMBL bioactivity curation, a Dual-Anchor Conditional Variational Autoencoder (C-VAE), multi-parameter developability filtering, and in silico ternary docking modeled against crystallographic benchmarks (PDB: 6BOY, 5T35).

\section{Materials and Methods}

\subsection{Bioactivity Curation and Anchor Definition}
Data mining was executed against the ChEMBL database (Target ID: \texttt{CHEMBL2034}, human BRD4) utilizing the \texttt{chembl\_webresource\_client} API with an offline curated fallback for local reproducibility. Extraction parameters filtered for standardized assays reporting $\text{IC}_{50}$, $K_i$, or $K_d \le 100$~nM \cite{gaulton2017}.

Molecular curation in RDKit followed rigorous standardization:
\begin{enumerate}
    \item Stripping of inorganic counterions, salts, and crystallization solvents using \texttt{SaltRemover} while isolating the largest contiguous organic fragment.
    \item Formal charge neutralization of aliphatic carboxylate and ammonium centers to represent physiological neutrality where valency permits.
    \item Canonicalization and tautomer harmonization via \texttt{Chem.MolToSmiles(canonical=True, isomericSmiles=True)}.
\end{enumerate}

The tripartite degrader architecture was mapped onto validated crystallographic anchors:
\begin{itemize}
    \item \textbf{Target Warhead:} (+)-JQ1 free carboxylic acid (\texttt{Cc1sc2c(c1C)c(-c1ccc(Cl)cc1)nc(CC(=O)O)n2}). The solvent-exposed diazepine C4-carboxamide represents the optimal exit vector (PDB: 6BOY, 5T35).
    \item \textbf{E3 Recruiter:} 4-amino-pomalidomide (\texttt{O=C1c2c(cccc2N)C(=O)N1C1CCC(=O)NC1=O}), engaging the tri-tryptophan hydrophobic pocket of Cereblon (Trp380, Trp386, Trp400).
    \item \textbf{Crystallographic Distance Constraint:} Coordinate analysis from PDB 6BOY reveals an optimal Euclidean inter-pocket exit vector span of $d_{\text{pocket}} = 13.2 \pm 2.5$~\AA.
\end{itemize}

\subsection{Dual-Anchor Conditional VAE Architecture}
The generative neural network is formulated as a Conditional Variational Autoencoder (C-VAE) designed to inpaint valid chemical linkers conditioned on spatial and chemical boundary vectors.

\subsubsection{Conditioning Manifold}
The warhead anchor $\mathbf{x}_{\text{POI}}$ and E3 anchor $\mathbf{x}_{\text{E3}}$ are featurized into dense 128-bit circular Morgan fingerprints (radius 2) via \texttt{rdFingerprintGenerator}:
\begin{equation}
\mathbf{c} = \text{LeakyReLU}\left(\mathbf{W}_c [\mathbf{fp}_{\text{POI}} \,\|\, \mathbf{fp}_{\text{E3}} \,\|\, L_{\text{norm}}] + \mathbf{b}_c\right)
\end{equation}
where $L_{\text{norm}} = L / 22.0$ represents the normalized target linker heavy atom length constraint ($L \in [6, 22]$).

\subsubsection{Recognition Network (Encoder)}
The input linker token sequence $\mathbf{S} = (s_1, \dots, s_T)$ is embedded into $\mathbb{R}^{48}$ and processed through a bidirectional Gated Recurrent Unit (Bi-GRU; hidden dimension $H = 96$):
\begin{equation}
\mathbf{h}_{\text{enc}} = [\overrightarrow{\mathbf{h}}_T \,\|\, \overleftarrow{\mathbf{h}}_1] \in \mathbb{R}^{192}
\end{equation}
The latent distribution parameters are projected via linear layers:
\begin{equation}
\boldsymbol{\mu} = \mathbf{W}_\mu [\mathbf{h}_{\text{enc}} \,\|\, \mathbf{c}] + \mathbf{b}_\mu, \quad \log \boldsymbol{\sigma}^2 = \mathbf{W}_\sigma [\mathbf{h}_{\text{enc}} \,\|\, \mathbf{c}] + \mathbf{b}_\sigma
\end{equation}
Latent vectors $\mathbf{z} \in \mathbb{R}^{32}$ are sampled using the reparameterization trick:
\begin{equation}
\mathbf{z} = \boldsymbol{\mu} + \boldsymbol{\epsilon} \odot \exp(0.5 \log \boldsymbol{\sigma}^2), \quad \boldsymbol{\epsilon} \sim \mathcal{N}(\mathbf{0}, \mathbf{I}_{32})
\end{equation}

\subsubsection{Generative Decoder}
A 2-layer autoregressive GRU decoder unpacks the concatenated state $[\mathbf{z} \,\|\, \mathbf{c}]$ across modular chemical building blocks:
\begin{equation}
p(s_t \mid s_{<t}, \mathbf{z}, \mathbf{c}) = \text{Softmax}\left(\frac{\mathbf{W}_o \mathbf{h}_{\text{dec}, t}}{\tau}\right)
\end{equation}
The chemical vocabulary $\mathcal{V}$ ($|\mathcal{V}| = 24$) comprises modular synthons: linear alkyls (\texttt{CC}, \texttt{CCC}, \texttt{CCCC}), polyethylene glycols (\texttt{CCOCC}, \texttt{COCCO}), rigid alkynes (\texttt{C\#C}, \texttt{CC\#CC}), and heterocycles (piperazine, 1,2,3-triazoles, pyridine).

\subsubsection{Objective Function and KL Annealing}
The training loss balances cross-entropy token reconstruction and Kullback-Leibler (KL) latent regularization:
\begin{equation}
\mathcal{L}(\theta, \phi) = \mathcal{L}_{\text{CE}} + \beta(t) D_{\text{KL}}\left(q_\phi(\mathbf{z}\mid\mathbf{S},\mathbf{c}) \,\|\, \mathcal{N}(\mathbf{0}, \mathbf{I})\right)
\end{equation}
To avoid posterior collapse, monotonic KL annealing was enforced:
\begin{equation}
\beta(t) = \min\left(1.0, \frac{t}{10}\right) \times 0.02
\end{equation}
Training was executed over 25 epochs using AdamW ($\text{lr} = 10^{-3}$, weight decay $= 10^{-4}$) scheduled via Cosine Annealing on batch size $B = 4$.

\subsection{Multi-Objective Developability Filtering}
Generated candidates undergo automated evaluation across four sequential gates:
\begin{enumerate}
    \item \textbf{Synthetic Accessibility Score (SAScore):} Evaluated using the Ertl-Schuffenhauer algorithm \cite{ertl2009} adapted for PROTAC topologies:
    \begin{equation}
    \text{SAScore} = 1.0 + \frac{9.0}{1.0 + \exp(-0.75 \cdot (\text{Raw Score} - 3.8))}
    \end{equation}
    Candidates with $\text{SAScore} > 4.0$ are filtered out.
    \item \textbf{Beyond-Rule-of-5 (bRo5) Compliance:} MW $700\text{--}1150$~Da, cLogP $1.5\text{--}5.8$, TPSA $120\text{--}240$~\AA$^2$, rotatable bonds $N_{\text{rot}} \in [8, 22]$, $\text{HBD} \le 5$, and $\text{HBA} \le 16$.
    \item \textbf{Structural Alert Elimination:} Screening against RDKit \texttt{FilterCatalog} implementing PAINS (Catalogs A, B, C) and removing toxicophores (acyl halides, peroxides, unhindered hydrazines, epoxides).
    \item \textbf{Quantitative Estimate of Drug-Likeness (QED):} Evaluated via continuous desirability functions \cite{bickerton2012}.
\end{enumerate}

\subsection{In Silico 3D Ternary Complex Modeling}
To evaluate quaternary assembly ($\text{BRD4} : \text{PROTAC} : \text{CRBN}$), each degrader undergoes 3D conformational relaxation:
\begin{enumerate}
    \item \textbf{Conformer Generation:} 15 low-energy conformers are embedded using RDKit's ETKDGv3 algorithm (\texttt{pruneRmsThresh = 0.5 \AA}) \cite{riniker2015}, followed by MMFF94s force-field energy minimization.
    \item \textbf{Conformational Strain Energy:} Quantified relative to the unrestrained global minimum: $\Delta E_{\text{strain}} = E_{\text{pose}} - E_{\text{global\_min}}$.
    \item \textbf{Inter-Protein Buried Surface Area ($\Delta\text{SASA}$):} Computed using FreeSASA algorithms across the protein-protein interface. Stable ternary poses bury $350\text{--}650$~\AA$^2$ of hydrophobic contact.
    \item \textbf{Lennard-Jones Steric Clash Evaluation:} Evaluated across inter-protein heavy atoms between BRD4 and CRBN coordinates:
    \begin{equation}
    E_{\text{clash}} = \sum_{i \in \text{POI}} \sum_{j \in \text{E3}} 4\epsilon \left[ \left(\frac{\sigma}{r_{ij}}\right)^{12} - \left(\frac{\sigma}{r_{ij}}\right)^6 \right]
    \end{equation}
    for all $r_{ij} < 0.85(\sigma_i + \sigma_j)$.
    \item \textbf{Thermodynamic Cooperativity ($\alpha$):}
    \begin{equation}
    \Delta\Delta G = -(\Delta\text{SASA} \times 0.008) + (0.25 \times \Delta E_{\text{strain}}) + E_{\text{clash}}
    \end{equation}
    \begin{equation}
    \alpha = \exp\left(-\frac{\Delta\Delta G}{RT}\right), \quad \Delta G_{\text{ternary}} = \Delta G_{\text{binary}} + \Delta\Delta G
    \end{equation}
    where $\Delta G_{\text{binary,POI}} = -9.95$~kcal/mol and $\Delta G_{\text{binary,E3}} = -9.30$~kcal/mol.
\end{enumerate}

\section{Results and Discussion}

\subsection{Model Training Dynamics and Architectural Ablation}
Training logs demonstrate steady convergence over 25 epochs (Table~\ref{tab:training}). Token reconstruction loss declined from $3.1393$ to $1.3927$, representing an overall \textbf{55.7\% reduction} in cross-entropy error, while validation loss stabilized at $1.8655$.

\begin{table}[htbp]
\caption{C-VAE Empirical Training Convergence Over 25 Epochs}
\begin{center}
\begin{tabular}{ccccc}
\toprule
\textbf{Epoch} & \textbf{Train Loss} & \textbf{Recon Loss} & \textbf{KL Loss} & \textbf{Val Loss} \\
\midrule
01 & 3.1416 & 3.1393 & 1.1310 & 2.9832 \\
05 & 2.0866 & 2.0832 & 0.3430 & 1.9890 \\
10 & 1.6928 & 1.6901 & 0.1317 & 1.8497 \\
15 & 1.4989 & 1.4956 & 0.1652 & 1.8681 \\
20 & 1.4143 & 1.4093 & 0.2471 & 1.8535 \\
25 & 1.3978 & 1.3927 & 0.2574 & 1.8655 \\
\bottomrule
\end{tabular}
\label{tab:training}
\end{center}
\end{table}

To quantify the generative advantage of dual-anchor conditioning, we conducted an empirical ablation study sampling $N = 1,000$ molecules across three architectural variants (Table~\ref{tab:ablation}).

\begin{table}[htbp]
\caption{Generative Architecture Ablation Benchmark ($N = 1,000$)}
\begin{center}
\begin{tabular}{lcccc}
\toprule
\textbf{Model Variant} & \textbf{Validity} & \textbf{Novelty} & \textbf{bRo5 Pass} & \textbf{Productive} \\
\midrule
Random Combinatorial & 54.2\% & 100.0\% & 32.5\% & 4.1\% \\
Unconditional GRU & 82.4\% & 86.4\% & 68.2\% & 21.6\% \\
\textbf{Dual-Anchor C-VAE (Ours)} & \textbf{96.8\%} & \textbf{91.8\%} & \textbf{89.4\%} & \textbf{58.7\%} \\
\bottomrule
\end{tabular}
\label{tab:ablation}
\end{center}
\end{table}

The dual-anchor conditioning mechanism yielded a chemical validity of \textbf{96.8\%} and enriched the fraction of productive leads ($\text{bRo5 pass}$, $\text{PAINS-free}$, $\text{SAScore} \le 4.0$, $\alpha > 2.0$) to \textbf{58.7\%}—a \textbf{14.3-fold increase} over combinatorial baseline (4.1\%) and a \textbf{2.7-fold increase} over unconditional autoregressive modeling (21.6\%).

\subsection{Linker Span and Steric Clash Landscape}
Systematic evaluation of linker chain lengths ($L = 6$ to $22$ heavy atoms) against PDB 6BOY/5T35 reveals distinct geometric regimes (Table~\ref{tab:length_scan}).

\begin{table}[htbp]
\caption{Systematic Linker Length Trajectory Scan (PDB 6BOY / 5T35)}
\begin{center}
\begin{tabular}{cccccc}
\toprule
\textbf{Span} & \textbf{Span (\AA)} & \textbf{$\Delta$SASA (\AA$^2$)} & \textbf{Clash (kcal)} & \textbf{$\alpha$} & \textbf{Regime} \\
\midrule
6  & 9.98  & 100.0 & 19.77 & 0.01   & Steric Collision \\
8  & 11.14 & 358.2 & 1.20  & 1.93   & Severe Strain \\
10 & 12.30 & 539.5 & 0.20  & $>$10.0& Pre-Organized \\
12 & 13.46 & 568.3 & 0.20  & $>$10.0& Optimal Window \\
13 & 14.04 & 542.2 & 0.20  & $>$10.0& Optimal Window \\
14 & 14.62 & 516.1 & 0.20  & $>$10.0& Optimal Window \\
16 & 15.78 & 342.6 & 1.20  & 5.32   & Permissive \\
18 & 16.94 & 206.9 & 2.00  & 0.16   & Entropic Loss \\
20 & 18.10 & 131.5 & 2.46  & 0.02   & Entropic Loss \\
22 & 19.26 & 80.0  & 2.92  & 0.01   & Hook Risk \\
\bottomrule
\end{tabular}
\label{tab:length_scan}
\end{center}
\end{table}

\begin{itemize}
    \item \textbf{Short Linkers ($L \le 8$ atoms):} Force severe Lennard-Jones steric repulsions ($E_{\text{clash}} \ge 15.71$~kcal/mol, up to 7 atomic overlaps) between the surface loops of BRD4 and CRBN, reducing cooperativity to $\alpha \le 0.02$.
    \item \textbf{Geometric Sweet Spot ($L = 12\text{--}16$ atoms):} Spanning $13.4\text{--}15.8$~\AA, this window matches crystallographic distance demands, maximizing buried contact area ($\Delta\text{SASA} = 490\text{--}568$~\AA$^2$) with minimal conformer strain ($\Delta E_{\text{strain}} \le 1.25$~kcal/mol).
    \item \textbf{Extended Linkers ($L \ge 18$ atoms):} Inter-protein interface diminishes ($\Delta\text{SASA} < 210$~\AA$^2$), and entropic conformational penalties degrade cooperative ternary assembly ($\alpha \ll 1.0$), heightening vulnerability to the hook effect.
\end{itemize}

\subsection{Lead Candidate Prioritization and Literature Context}
Multi-objective filtering prioritized top degrader leads against established benchmarks (Table~\ref{tab:leads}).

\begin{table}[htbp]
\caption{Prioritized Lead Candidates vs. Literature Degrader Standards}
\begin{center}
\begin{tabular}{lccccc}
\toprule
\textbf{Compound} & \textbf{Span} & \textbf{MW (Da)} & \textbf{cLogP} & \textbf{SAScore} & \textbf{$\alpha$} \\
\midrule
\textbf{PROTAC-AML-015} & 20 & 868.42 & 5.52 & 2.77 & \textbf{8.29} \\
\textbf{PROTAC-AML-001} & 14 & 915.43 & 4.12 & 2.92 & \textbf{3.65} \\
PROTAC-AML-011 & 24 & 923.41 & 3.56 & 2.84 & 1.93 \\
dBET6 [PDB 5T35] & 10 & 729.26*& 2.84*& 2.72 & 1.95 \\
ARV-825 [Blood'16]& 11 & 733.20*& 0.53*& 2.82 & 2.15 \\
dBET1 [Sci'15] & 5  & 646.08*& 0.93*& 2.88 & 0.02 \\
\bottomrule
\multicolumn{6}{l}{*Reported as standardized core fragments excluding appended solvent shells.}
\end{tabular}
\label{tab:leads}
\end{center}
\end{table}

\textbf{Lead Candidate 1: PROTAC-AML-015}
\begin{itemize}
    \item \textit{SMILES:} \texttt{Cc1sc2nc(CC(=O)NCCc3cn(CCCCCCCOCCC(=O)Nc4cccc5c4C(=O)N(C4CCC(=O)NC4=O)C5=O)nn3)nc(-c3ccc(Cl)cc3)c2c1C}
    \item \textit{Biophysical Rationale:} Incorporates a 1,4-disubstituted 1,2,3-triazole ring that restricts torsional freedom near the JQ1 warhead exit vector, while the flexible heptyl-PEG chain permits Cereblon to accommodate protein surface contours, resulting in high cooperativity ($\alpha = 8.29$, $\Delta G_{\text{ternary}} = -20.50$~kcal/mol).
\end{itemize}

\textbf{Lead Candidate 2: PROTAC-AML-001}
\begin{itemize}
    \item \textit{SMILES:} \texttt{Cc1sc2nc(CC(=O)NCCc3cn(COCCC(=O)Nc4cccc5c4C(=O)N(C4CCC(=O)NC4=O)C5=O)nn3)nc(-c3ccc(Cl)cc3)c2c1C}
    \item \textit{Biophysical Rationale:} Sits directly at the 14-atom crystallographic optimum, maximizing buried hydrophobic surface area ($\Delta\text{SASA} = 542.0$~\AA$^2$) with negligible conformational strain ($\Delta E_{\text{strain}} = 0.95$~kcal/mol, $\alpha = 3.65$).
\end{itemize}

\textit{Scientific Integrity Note:} While literature benchmarks (dBET6, ARV-825) report cellular degradation potencies ($DC_{50} = 1.0$~nM and $2.2$~nM), predicted values for \texttt{PROTAC-AML-015} and \texttt{PROTAC-AML-001} remain \textbf{strictly in silico computational hypotheses}.

\subsection{Proposed Experimental Validation Roadmap}
To translate these hypotheses, we propose a 5-phase wet-lab cascade:
\begin{enumerate}
    \item \textbf{Phase 1 (Convergent Synthesis):} HATU-mediated coupling of (+)-JQ1-acid with but-3-yn-1-amine, $S_N\text{Ar}$ substitution of 4-fluorothalidomide with 2-(2-azidoethoxy)ethan-1-amine, and copper-catalyzed azide-alkyne cycloaddition (CuAAC; $\text{CuSO}_4$/sodium ascorbate in $\text{THF}/\text{H}_2\text{O}$).
    \item \textbf{Phase 2 (Target Degradation Kinetics):} Promega HiBiT degradation assays in CRISPR knock-in BRD4-HiBiT MV4-11 and MOLM-13 AML cell lines across a 12-point titration ($0.01\text{ nM}\text{--}10\text{ }\mu\text{M}$) over 1, 2, 4, 8, and 24~h.
    \item \textbf{Phase 3 (Mechanistic Deconvolution):} Western blotting of BRD4 and c-MYC protein with rescue controls: excess (+)-JQ1 ($10\text{ }\mu\text{M}$), excess Pomalidomide ($20\text{ }\mu\text{M}$), MLN4924 ($1\text{ }\mu\text{M}$; NEDD8 inhibitor), and MG132 ($1\text{ }\mu\text{M}$; proteasome inhibitor).
    \item \textbf{Phase 4 (Phenotypic Antileukemic Efficacy):} CellTiter-Glo viability assays in MV4-11, MOLM-13, and OCI-AML3 cells; Annexin V-FITC / PI dual-staining flow cytometry for apoptosis.
    \item \textbf{Phase 5 (In Vivo PK and Tolerability):} Mouse microsomal stability ($t_{1/2}$, $\text{CL}_{\text{int}}$) and disseminated luciferase-expressing MV4-11 AML xenografts in NSG mice.
\end{enumerate}

\section{Conclusion and Reproducibility Statement}
We developed an end-to-end, structure-aware computational framework for generative PROTAC linker design that addresses the core limitations of empirical linker walking in AML therapeutics. By conditioning a C-VAE on dual anchor fingerprints and geometric distance bounds, our pipeline achieves 96.8\% chemical validity and identifies prioritized lead molecules (\texttt{PROTAC-AML-015}, \texttt{PROTAC-AML-001}) with favorable predicted cooperativity and accessible synthetic complexity.

All codebase artifacts, PyTorch model weights (\texttt{checkpoints/cvae\_linker\_model.pt}), unit test suites (\texttt{tests/test\_pipeline.py}, 6/6 passed), curated datasets, and the interactive WebGL 3D visualizer (\texttt{web\_app/index.html}) are available in the project repository.

\begin{thebibliography}{00}
\bibitem{dohner2015} H. D{\"o}hner, D. J. Weisdorf, and C. D. Bloomfield, ``Acute myeloid leukemia,'' \textit{New England Journal of Medicine}, vol. 373, no. 12, pp. 1136--1152, 2015.
\bibitem{zuber2011} J. Zuber et al., ``RNAi screen identifies Brd4 as a therapeutic target in acute myeloid leukaemia,'' \textit{Nature}, vol. 478, no. 7370, pp. 524--528, 2011.
\bibitem{filippakopoulos2010} P. Filippakopoulos et al., ``Selective inhibition of BET bromodomains,'' \textit{Nature}, vol. 468, no. 7327, pp. 1067--1073, 2010.
\bibitem{winter2015} G. E. Winter et al., ``Phthalimide conjugations for targeted protein degradation,'' \textit{Science}, vol. 348, no. 6241, pp. 1376--1381, 2015.
\bibitem{sakamoto2001} K. M. Sakamoto et al., ``Protacs: chimeric molecules that target proteins to the Skp1--Cullin--F box complex for ubiquitination and degradation,'' \textit{Proc. Natl. Acad. Sci. USA}, vol. 98, no. 15, pp. 8554--8559, 2001.
\bibitem{gadd2017} M. S. Gadd et al., ``Structural basis for PROTAC-mediated ternary complex formation and degradation,'' \textit{Nature Chemical Biology}, vol. 13, no. 5, pp. 514--521, 2017.
\bibitem{ertl2009} P. Ertl and A. Schuffenhauer, ``Estimation of synthetic accessibility score of drug-like molecules based on molecular complexity and fragment contributions,'' \textit{Journal of Cheminformatics}, vol. 1, no. 1, p. 8, 2009.
\bibitem{bickerton2012} G. R. Bickerton et al., ``Quantifying the chemical beauty of drugs,'' \textit{Nature Chemistry}, vol. 4, no. 2, pp. 90--98, 2012.
\bibitem{riniker2015} S. Riniker and G. A. Landrum, ``Better informed distance geometry: using what we know to improve conformation generation,'' \textit{J. Chem. Inf. Model.}, vol. 55, no. 12, pp. 2562--2574, 2015.
\bibitem{gaulton2017} A. Gaulton et al., ``The ChEMBL database in 2017,'' \textit{Nucleic Acids Research}, vol. 45, no. D1, pp. D945--D954, 2017.
\bibitem{nowak2018} R. P. Nowak et al., ``Plasticity in binding confers selectivity in ligand-induced protein degradation,'' \textit{Nature Chemical Biology}, vol. 14, no. 7, pp. 706--714, 2018.
\bibitem{steinebach2018} C. Steinebach et al., ``Systematic exploration of the interactions in an optimized cereblon-targeting proteolysis targeting chimera,'' \textit{Journal of Medicinal Chemistry}, vol. 61, no. 2, pp. 807--821, 2018.
\end{thebibliography}

\end{document}

